% SPDX-FileCopyrightText: © 2026 Idiap Research Institute <contact@idiap.ch>
%
% SPDX-FileContributor: Andre Freitas <andre.freitas@idiap.ch>
% Neuro-symbolic AI Group
%
% SPDX-FileContributor: Danilo Gusicuma <danilo.gusicuma@idiap.ch>
% Neuro-symbolic AI Group
%
% SPDX-License-Identifier: MIT
% Generated from step-explanations.json
\paragraph{Step-by-step reconstruction.}
A state-convergence argument uses a norm that distinguishes gauge representatives even when they represent the same observable state.
The following is a source-guided illustrative reconstruction. Internal procedures are prescriptions; the native declarations do not establish their execution.
\begin{enumerate}[leftmargin=*,label=\arabic*.]
\item \textbf{Locate the representational obstruction} (\texttt{obstruction\_\allowbreak{}pressure}).
\begin{description}[style=standard,labelindent=0pt,leftmargin=5.5em,labelwidth=5em,labelsep=0.5em,align=left,font=\normalfont\bfseries,itemsep=0.15em,parsep=0pt,topsep=0.3em]
\item[Purpose] Distinguish conceptual obstruction from unfinished calculation.
\item[Input] Failed argument; norm’s identity distinction
\item[Method] Trace the failed implication to the role played by the carrier, invariant or equivalence relation.
\item[Check] An obstruction diagnosis must explain why existing objects fail at the named boundary.
\item[Output] Obstruction diagnosis: representative-dependent convergence criterion
\item[Next] Mediation consumes this diagnosis; a merely missing calculation should remain a derivation task.
\item[Example] Two gauge-equivalent representatives can be far apart in the chosen norm, so representative convergence need not express physical-state convergence.
\end{description}
\item \textbf{Propose the mediating object} (\texttt{mediation\_\allowbreak{}hypothesis}).
\begin{description}[style=standard,labelindent=0pt,leftmargin=5.5em,labelwidth=5em,labelsep=0.5em,align=left,font=\normalfont\bfseries,itemsep=0.15em,parsep=0pt,topsep=0.3em]
\item[Purpose] Propose a carrier for the intended question.
\item[Input] Representative-dependent identity and convergence obstruction
\item[Method] Propose a quotient or other mediator, naming what it would preserve and its first possible failure regime.
\item[Check] The desired convergence theorem cannot be hidden inside the mediator's definition.
\item[Output] Mediation hypothesis: gauge-equivalence classes
\item[Next] Transfer must turn this proposal into specified maps; a name alone is insufficient.
\item[Example] Replace the proposed carrier of raw representatives by classes under gauge equivalence, while leaving existence of suitable dynamics as an obligation.
\end{description}
\item \textbf{Specify the bridge} (\texttt{conceptual\_\allowbreak{}bridge}).
\begin{description}[style=standard,labelindent=0pt,leftmargin=5.5em,labelwidth=5em,labelsep=0.5em,align=left,font=\normalfont\bfseries,itemsep=0.15em,parsep=0pt,topsep=0.3em]
\item[Purpose] Specify the mediator’s maps and compatibility burden.
\item[Input] Quotient proposal; intended observable interpretation
\item[Method] Specify projection from representatives to classes and the compatibility needed for dynamics to descend.
\item[Check] Equivalent representatives must induce the same class evolution for the bridge to be well defined.
\item[Output] Projection bridge; open dynamics-compatibility obligation
\item[Next] Deduction uses this bridge conditionally; an absent compatibility law blocks unconditional transport.
\item[Example] For representative evolution F, require the class of F(x) to equal the class of F(y) whenever x and y are gauge equivalent.
\end{description}
\item \textbf{Derive a bridge consequence} (\texttt{bridge\_\allowbreak{}consequence}).
\begin{description}[style=standard,labelindent=0pt,leftmargin=5.5em,labelwidth=5em,labelsep=0.5em,align=left,font=\normalfont\bfseries,itemsep=0.15em,parsep=0pt,topsep=0.3em]
\item[Purpose] State the bridge’s conditional consequence.
\item[Input] Mediating bridge; open compatibility assumptions
\item[Method] Formulate the transported statement on the new carrier and retain each bridge premise.
\item[Check] The consequence must concern equivalence classes without silently restoring representative identity.
\item[Output] Class-level evolution or convergence claim, conditional on its premises
\item[Next] Refutation tests this consequence; convergence estimates remain separate from well-definedness.
\item[Example] A quotient trajectory can be defined from F only if changing the representative leaves its next class unchanged; convergence would then be a class-level question.
\end{description}
\item \textbf{Test preservation} (\texttt{refutation\_\allowbreak{}assessment}).
\begin{description}[style=standard,labelindent=0pt,leftmargin=5.5em,labelwidth=5em,labelsep=0.5em,align=left,font=\normalfont\bfseries,itemsep=0.15em,parsep=0pt,topsep=0.3em]
\item[Purpose] Test whether the representation is well defined.
\item[Input] Bridge consequence; compatibility premise
\item[Method] Search for equivalent inputs whose evolved outputs cease to be equivalent; localize the precise preservation law at fault.
\item[Check] A bridge fails this test when one physical input class receives two distinct output classes.
\item[Output] Refutation assessment: well-definedness defect or bounded surviving claim
\item[Next] Pass the finding to routing with the retained bridge; no found defect is not a proof.
\item[Example] Choose equivalent x and y with inequivalent F(x) and F(y). The proposed class dynamics is ambiguous, independently of any convergence estimate.
\end{description}
\item \textbf{Route the conceptual repair} (\texttt{refinement\_\allowbreak{}routing}).
\begin{description}[style=standard,labelindent=0pt,leftmargin=5.5em,labelwidth=5em,labelsep=0.5em,align=left,font=\normalfont\bfseries,itemsep=0.15em,parsep=0pt,topsep=0.3em]
\item[Purpose] Choose proof work or bridge reconstruction.
\item[Input] Conceptual bridge + its refutation assessment
\item[Method] Compare target contribution, validation strength and the exact preservation defect; distinguish an informative repair from a repeated blocked route.
\item[Check] A work-ready bridge needs explicit obligations that address the bottleneck.
\item[Output] Refinement assessment: compatibility-repair obligation
\item[Next] Original policy routes work or release and stops refutation; native routing is a returned value.
\item[Example] Keep quotienting as a candidate representation but require a compatible evolution law; do not advertise the quotient name as a solved convergence problem.
\end{description}
\end{enumerate}
\paragraph{Changed-input contrast.}
If equivalent representatives evolve into equivalent classes, the proposed quotient passes this well-definedness test. Convergence still requires its own argument; compatibility alone does not establish it.
